\documentclass{jsarticle}
\usepackage{amsmath}
\usepackage{amssymb}
\begin{document}
\title{}
\author{}
\date{}
\maketitle


    $$ f(R) = - R_{ij} $$
    $$ f(L) = R_{ij} $$
    ベクトルでスカラーに向きを加えている。
    右回転のときに反重力
    $$ f(L) - f(R) = -2 R_{ij} $$
    $$i {d \over dt}g_{ij} = - 2 R_{ij} $$
    質量差エネルギーとエネルギー密度で、
    一般相対性理論の多様体積分と特殊相対性理論の光速度不変式を使い、
    反重力発生器を作る原理を、大域的積分多様体のガンマ関数で、
    各方程式を述べている。
    宇宙人を真剣に述べているのは、懐疑的に感じたが、
    全体的に、説得力のある本自体になっている。　

\end{document}
